\chapter{Introduction}\label{ch:intro}

\section{The question}
When two ways of holding the same risk trade at different prices, an arbitrageur can buy the cheap one, sell the dear one
and wait for the prices to meet. If arbitrage were free and unlimited, such gaps would never be large or long-lived.
Shleifer and Vishny (1997) and Gromb and Vayanos (2010) explain why it is neither: arbitrage needs capital, it carries
risk while the gap persists, it costs money to trade, and it often needs a short position that may be costly or not
allowed. Each of these limits lets prices drift away from each other.

This thesis asks how strongly these limits bind in the Indian equity market, which of them matter, and how much
mispricing they explain. India is a good place to ask. Almost every limit changes on known dates: the securities
transaction tax (STT) on futures was raised three times between 2023 and 2026; only stocks in the futures and options
(F\&O) segment can be sold short or shorted through stock futures, and stocks enter and leave that segment every few
months; the assets of arbitrage mutual funds grew several times after a change in the taxation of debt funds in 2023; and
position limits are set by formula and were changed in October 2025.

India also has an unusual rule. When market-wide open interest in a stock's derivatives exceeds 95\% of its market-wide
position limit (MWPL), the stock enters an \emph{F\&O ban}: every participant, not only arbitrageurs, may only trade to
reduce positions until open interest falls to 80\%. Between 2015 and October 2026 this happened 1,788 times to 141 stocks.
To my knowledge no study has estimated its effects. The ban is the sharpest lens on one limit, but it is one lens among
several; the thesis studies the limits side by side.

\section{Approach}
The work followed a fixed order, designed to prevent fitting the theory and the data to each other.
\begin{enumerate}
\item \emph{Theory first.} One model covers all the limits; each enters as a parameter. The ban is modelled as the rule
actually works. Every hypothesis is derived from the model with the result that would contradict it (Chapter~\ref{ch:theory}).
\item \emph{A frozen test plan.} Samples, variables, specifications and falsifiers were written down and frozen with a
time-stamped record before any outcome was related to any treatment. Every detail the plan left open was fixed in a dated
log before the test that needed it was run (Appendix~\ref{app:plan}).
\item \emph{Data from primary sources.} All data come from public exchange archives and are built with scripts that check
every step against the exchange's own rules (Chapter~\ref{ch:data}).
\item \emph{Estimators tested on simulated data} from the model before use.
\item \emph{All results reported,} including those that contradict the model.
\end{enumerate}

\section{Main findings}
\begin{enumerate}
\item \textbf{Bans bind on trading.} In a ban, futures trading falls by about two log points and the delivery share of cash
trading rises. When the rule was enforced on peak intraday positions (March 2020 to September 2025), the change in open
interest equals minus the volume: every futures trade closed a long and a short. Under end-of-day enforcement before March
2020 this identity did not hold.
\item \textbf{Correction is not one-sided as predicted.} The model predicts that rich futures stop being corrected in a ban
because nobody may open new cash-and-carry positions. The data show the opposite: rich futures revert faster in a ban, and
cheap futures more slowly. The effect is concentrated in the first days of a spell and survives a correction for noise in
daily prices.
\item \textbf{The basis falls when a ban starts.} Against stocks that came close to the trigger but did not cross it, the
basis is about 6 basis points (bp) lower in the first days of a ban, more when open interest had been building fast, and
it returns to its earlier level after about a week. Utilisation does not bunch below the trigger, and as utilisation rises
the basis is more often rich.
\item \textbf{Costs, capital and the band.} The basis follows a wide, asymmetric band, about $-35$ to $+24$\bp\ around each
stock's median. The upper part of the distribution rose around the 2024 and 2026 STT increases, but by more than the cost
changes and mostly as a common shift. A new cash-and-carry round trip costs 27 to 31\bp, which places it in the outer tail
of the basis; rolling an existing position costs 4 to 8\bp. Inflows into arbitrage funds are followed by a lower upper
tail, but the evidence is weak.
\item \textbf{Short-sale access and the cross-section.} After a stock enters the F\&O segment its price responds faster to
market-wide news; drift after bad news does not weaken. A constraint index built only from frictions, for every listed
stock, predicts slower price adjustment, also after controlling for size; its link with drift after bad news is a size
effect.
\end{enumerate}

\section{Contribution}
The thesis studies the F\&O ban for the first time; models it as a reduce-only rule for every participant and shows that
this makes correction one-sided in theory and the level effect ambiguous in sign; shows that changes in the position limit
cannot serve as an instrument for ban status; estimates the no-arbitrage band of Indian stock futures over a long daily
sample with threshold and smooth-transition models; and publishes a constraint index for all Indian stocks. The closest
earlier work on Indian stock futures estimates bands from intraday data over short samples without the ban (Aggarwal,
2015; Shankar, Sankar and Kiran Kumar, 2019); Saurav, Srivastava and Agarwalla (2026) study the basis around earnings
announcements and short-sale constraints.

\section{Outline}
Chapter~\ref{ch:inst} describes the Indian rules. Chapter~\ref{ch:theory} presents the model. Chapter~\ref{ch:data}
describes the data, the measures and the constraint index. Chapter~\ref{ch:band} estimates the band and relates it to
costs, capital and position limits. Chapter~\ref{ch:ban} studies the ban. Chapter~\ref{ch:elig} studies entry into and exit
from the F\&O segment and the cross-section of stocks. Chapter~\ref{ch:concl} concludes and discusses policy.
